% \chapter{Título capítulo: aprender a trabajar en LaTeX  ==> eliminar}

% \todo[inline]{Breves notas sobre como usar LaTeX  $\rightarrow$ \textbf{eliminar en la versión final} }

% Esta es una referencia a un artículo \cite{article-minimal}.

% Esta es una referencia a un capítulo dentro de un libro \cite{inbook-minimal}.

% Esta es una referencia a un libro \cite{book-minimal}.

% Esta es una referencia a un artículo dentro de los proceedings de un congreso \cite{inproceedings-full}.

% Esta es una referencia a una url \cite{misc-url}.

% \textbf{Esto está escrito en negrita}


% \emph{Esto está en enfatizado}

% \begin{center}
% Este texto está centrado
% \end{center}


% Esto es una lista:
% \begin{itemize}
% \item Primer elemento
% \item Segundo elemento
% \item Lista dentro de otra lista:
% 	\begin{itemize}
% 		\item Primer subelemento
% 		\item Segundo subelemento
% 	\end{itemize}
% \end{itemize}


% \section{Título sección}

% Esto es una descripción:
% \begin{description}
% \item[Palabra] descripción de la palabra
% \item[Palabra] descripción de la palabra
% \end{description}


% Y esto una lista numerada:
% \begin{enumerate}
% \item Elemento
% \item Elemento
% \item Elemento
% \item Elemento
% \end{enumerate}


% Podemos incluir una figura y referenciarla de esta forma \ref{fig:logo}. Además podemos poner la página en la que está: \pageref{fig:logo}

% \begin{figure}[htp]
% \begin{center}
% \includegraphics[scale=0.2]{figures/logo.eps}
% \caption{Esta es la etiqueta de la figura}
% \label{fig:logo}
% \end{center}
% \end{figure}


% \begin{figure}[htp]
% \begin{center}
% \framebox{\includegraphics[scale=0.2]{figures/logo.eps}}
% \caption{Esta es la etiqueta de la figura con borde}
% \label{fig:logo2}
% \end{center}
% \end{figure}


% \subsection{Título subsección}

% Existen varias formas de incluir ecuaciones matemáticas. Las más utilizadas son las siguientes:
% \begin{itemize}
%  \item Para introducir expresiones matemáticas en el texto, se utiliza como delimitador el símbolo del dólar. Por ejemplo $a \rightarrow b$.
%  \item Para introducir ecuaciones matemáticas, se utiliza el entorno \texttt{equation}:
% \begin{equation}
%  \gamma = \frac{\overline{\alpha}}{\sqrt{\beta}}
% \label{eq:equation_example_1}
% \end{equation}

% \begin{equation}
% E(v) =  \int^1_0 \int^1_0 \int^1_0 E_{int}(v(r,s,t)) + E_{ext}(v(r,s,t))drdsdt 
% \label{eq:equation_example_2}
% \end{equation}	
% Un subíndice se especifica con el guión bajo y un superíndice con el circunflejo. Si el super/subíndice contiene varios caracteres, estos deben estar delimitados por llaves. Consultar el manual para comprobar como se pueden introducir símbolos y expresiones matemáticas en latex.
% \end{itemize}



% \subsection{Título subsección}

% Esto es una subsección\footnote{Así se hace una nota a pie de página}.


% Así introducimos texto sin ningún tipo de formato latex:
% \begin{verbatim}
%   4 drwxr-xr-x  2 noelia imagen   4096 2005-09-12 12:09 figures
%   4 -rwxr--r--  1 noelia imagen    585 2005-09-12 16:56 Makefile
%   4 -rw-r--r--  1 noelia imagen    647 2005-09-12 17:38 memoria.aux
%   4 -rw-r--r--  1 noelia imagen   1011 2005-09-12 17:18 memoria.bbl
%   4 -rw-r--r--  1 noelia imagen   1171 2005-09-12 17:18 memoria.blg
%  16 -rw-r--r--  1 noelia imagen  13440 2005-09-12 17:38 memoria.dvi
%   4 -rw-r--r--  1 noelia imagen    412 2005-09-12 17:38 memoria.glg
%   4 -rw-r--r--  1 noelia imagen    188 2005-09-12 17:38 memoria.glo
%   4 -rw-r--r--  1 noelia imagen    241 2005-09-12 17:38 memoria.gls
%   4 -rw-r--r--  1 noelia imagen    299 2005-09-12 17:38 memoria.ist
%   4 -rw-r--r--  1 noelia imagen    283 2005-09-12 17:38 memoria.lof
%  12 -rw-r--r--  1 noelia imagen   8997 2005-09-12 17:38 memoria.log
%   4 -rw-r--r--  1 noelia imagen     87 2005-09-12 17:38 memoria.lot
% 108 -rw-r--r--  1 noelia imagen 103492 2005-09-12 17:18 memoria.pdf
% 280 -rw-r--r--  1 noelia imagen 278687 2005-09-12 17:18 memoria.ps
%   4 -rwxr--r--  1 noelia imagen   1887 2005-09-12 17:40 memoria.tex
%   4 -rw-r--r--  1 noelia imagen    735 2005-09-12 17:38 memoria.toc
% \end{verbatim}



% \subsubsection{Título subsubsección}

% La figura \ref{tab:tabla_ejemplo} muestra un ejemplo de tabla básica. Para crear una tabla se utiliza el entorno \texttt{tabular} dentro del flotante \texttt{table}. 

% %% PARA ENTENDER ESTE TEXTO, MEJOR COMPILAR Y LEER EL DVI/PS/PDF
% %% \ es un caracter reservado de latex por lo que para poder utilizarlo cuando escribimos debemos utilizar el comando \textbackslash
% %% &, {, } también son caracteres reservados. En este caso, para escribirlos en el texto, anteponemos \ al caracter

% Tras \texttt{\textbackslash begin\{tabular\}} se declara el número de columnas de la tabla. Cada columna se especifica con una letra:
% \begin{description}
%  \item [c] Si el texto en la columna está centrado
%  \item [l] Si el texto en la columna está alineado a la izquierda
%  \item [r] Si el texto en la columna está alineado a la derecha
% \end{description}

% \begin{table}[htp]
% \caption{Tabla de ejemplo}
% \label{tab:tabla_ejemplo}
% \begin{center}
%  \begin{tabular}{rlc}\hline
%   Fila 1 Col. 1 larálala & Fila 1 Col. 2 lara & Fila 1 Col. 3  lalalalalalalala\\ \hline
%   Fila 2 Col. 1 & Fila 2 Col. 2 & Fila 2 Col. 3 \\ \hline
%  \end{tabular}
% \end{center}
% \end{table}

% Así \texttt{\textbackslash begin\{tabular\}\{rrrr\}} indica 4 columnas con texto alineado a la derecha y \texttt{\textbackslash begin\{tabular\}\{cl\}} indica dos columnas, la primera centrada y la segunda alineada a la izquierda.

% Entre \texttt{\textbackslash begin\{tabular\}\{\ldots\}} y \texttt{\textbackslash end\{tabular\}} se escribe el contenido de la tabla. Las columnas se separan con \& y las filas con \textbackslash \textbackslash. Un ejemplo de una fila con tres columnas sería el siguiente:
% \begin{center}
%  \texttt{ aaa \& bbbb \& cccc \textbackslash \textbackslash}\footnote{Importante: debe coincidir el número de columnas en la declaración con el número de columnas que se escriben dentro de la tabla}. 
% \end{center}

% Para incluir lineas horizontales existe el comando \texttt{\textbackslash hline}. Tras \texttt{\textbackslash begin\{tabular\}\{\ldots\}} podemos incluir un \texttt{\textbackslash hline}, pero \textbf{ojo}, en cada fila \texttt{\textbackslash hline} siempre debe ir tras \textbackslash \textbackslash.

% También es posible incluir líneas verticales, pero los manuales de estilo no lo aconsejan. Las líneas verticales se definen en la declaración con barras verticales. Por ejemplo, con \{\textbar c  \textbar c \textbar c\}  se crearían 3 líneas verticales, una antes de la primera columna, otra entre la primera y la segunda columna, y la tercera, entre la segunda y tercera columna.

% Dentro de las tablas podemos incluir expresiones matemáticas, texto enfatizado, negrita, etc. Es posible también fusionar varias filas o varias columnas dentro de una misma tabla. Para ello, consultar en manuales los comandos/paquetes \texttt{multirow} y \texttt{multicol}.


% Nueva referencia \cite{milibro}.

